% SOURCE docs/release/2026-10-29/source/docs/writing-repair/2026-10-29/src/bb525bc600530286.tex CHAR 0
\chapter{Longitudinal treatment, monitoring and mediation}
\label{sw:longitudinal}

Policy choice also determines what the bank observes. Monitoring can change
record availability, and a recorded engagement measure may become a mediator
rather than a baseline control. The compact example here connects these roles
to the detailed longitudinal and mediation comparisons, so that an interface
preserves the question before routing it to an R estimator.

\section{The history is part of the intervention}

The servicing problem is naturally sequential. An initial action \(A_0\), such
as an early contact, can change a later financial condition \(L_1\). A second
action \(A_1\) is then chosen using that condition, and a later account cost
\(Y\) is observed. The variable \(L_1\) can be both a consequence of the first
action and a determinant of the second. Treating it as an ordinary baseline
control erases that temporal structure. Treating both actions as one binary
indicator can erase the policy being evaluated.

Consider a finite project mechanism in which \(A_0,L_1,A_1\) are binary and
\[
 \Pr(L_1=1\mid\operatorname{do}(A_0=0))=0.8,
 \qquad
 \Pr(L_1=1\mid\operatorname{do}(A_0=1))=0.4.
\]
Conditional mean cost, in CNY per account over a fixed horizon, is specified as
\[
 m(a_0,l,a_1)=100l-50a_1l+20a_0+10a_1.
\]
The initial action costs 20 and reduces the chance of the adverse condition.
The later action costs 10 and reduces adverse-condition cost by 50 when that
condition is present. These numbers define the example; they are not estimates
of actual bank effects.

For a second-stage rule \(d\), replacing the two action equations gives
\begin{equation}
 V(a_0,d)=\sum_{l=0}^{1}m(a_0,l,d(l))
             \Pr(L_1=l\mid\operatorname{do}(A_0=a_0)).
 \label{sw:finite-policy}
\end{equation}
Equation~\eqref{sw:finite-policy} follows by enumerating the two possible
conditions under the specified intervention. It yields
\[
 V(0,0)=80,\qquad V(1,l\mapsto l)=44,\qquad V(1,1)=50.
\]
The adaptive policy is better than always giving the later action because it
avoids paying for that action in accounts without the adverse condition.
No policy-learning or global-optimality theorem is needed for this comparison
of three explicitly specified rules. A proposed general optimizer would have
additional obligations.

Now hold \(A_0=0\) and consider an observed logger that gives \(A_1=1\) with
probability 0.9 when \(L_1=1\), and 0.1 otherwise. The observed treated group
contains a disproportionate share of accounts with the adverse condition.
Direct enumeration gives
\[
 E[Y\mid A_1=1,A_0=0]=58.64865\ldots,
 \qquad E[Y\mid A_1=0,A_0=0]=30.76923\ldots.
\]
The observational difference is approximately \(+27.88\). Yet the intervention
contrast that treats everyone at the second stage, at \(A_0=0\), is
\(V(0,1)-V(0,0)=50-80=-30\). The sign reversal is selection by an observed
time-varying condition in this mechanism. It shows why an interface must retain
the treatment history and the distribution over which the policy is valued.

In an application, evaluating an expression like
\eqref{sw:finite-policy} from observations requires substantive assumptions
connecting observed conditional distributions to the intervention model. It
also requires that relevant histories and actions have support. The example
supplies its structural law directly; it does not demonstrate that a bank's
logger, monitoring or missingness mechanisms satisfy those assumptions.

\section{Point treatment and sequential software have different contracts}

The selected tmle interface concerns binary point-treatment effects with
missingness and subject-ID options. It is a candidate for a suitable static
effect question; that does not make it an arbitrary longitudinal policy engine.
The versioned documentation also contains a treatment-missingness wording
difference between DESCRIPTION and the selected Delta argument.
Section~\ref{exec20:rfx-scope} resolves the selected implementation's contract:
missing treatment and baseline covariates are rejected, while the no-variation
route estimates an outcome mean. Section~\ref{sec:exec24-tmle-reuse} retains
the additional interval and resampling restrictions. A convenient argument
name cannot extend those admitted meanings. \cite{v12_sw_tmle}

stremr's selected sequential g-computation and TMLE documentation is closer to
the servicing problem. It accepts time points, sequential outcome-model
formulas, treatment intervention probabilities, monitoring intervention
probabilities and related weighting choices. Treatment and monitoring are
separate intervention components. A bank that increases contact may also
increase the chance that financial stress is recorded; holding monitoring
fixed and changing it are different policy questions. The interface should
name which of these is being evaluated. \cite{v12_sw_stremr}

The same documentation makes several operational qualifications explicit.
The documentation labels \texttt{Qstratify} nonfunctional and an error source,
while the source review in Section~\ref{sec:exec23-longr-workers} finds a
corresponding branch. That discrepancy remains a qualification question;
it is not evidence of failure in every native call. The proposed profile
withholds that option pending qualification. Weight truncation and stabilization
are choices, and the reviewed stopping loops can differ from the requested
cap. A wrapper should report the effective weights, actual updates, final
residuals and warnings after execution. A truncated weight or stopped numerical
iteration is not evidence that positivity or statistical accuracy has been
established. It is part of the approximation and diagnostic record.

survtmle addresses another closely related but distinct problem: right
censoring, competing failures and cumulative incidence. Its selected interface
supports two treatment values. A competing failure must not be recoded as
ordinary censoring simply to fit a common input schema. In a banking setting,
default, repayment, restructuring and account closure can correspond to
different outcome or observation events, depending on the question. The
adapter must preserve the cause, horizon and target, and the analyst must
justify their interpretation. \cite{v12_sw_survtmle}

The initial software decision is to evaluate separate R operations for static
TMLE, sequential g-computation/TMLE, and competing-risk cumulative incidence.
One common transport envelope can contain their inputs and outputs, but it
should not imply scientific substitutability. The proposed tests include
treatment-confounder feedback, monitoring-dependent recording, administrative
censoring, absent action support, and the distinction between a failure event
and a missing observation. These trials must assess the specific estimand and
sampling law; a successful fit on complete independent records does not
qualify a dynamic bank panel.

The age of a saved distribution is recorded as evidence about that distribution,
not as a verdict on project maintenance. In particular, older stremr and
survtmle snapshots do not prove abandonment. Conversely, a recent version
string does not demonstrate correctness of every documented feature. Build
compatibility and method validity require their own observations.

\section{Mediation changes the target again}

A servicing team may want to decompose an action's effect into a part operating
through customer engagement and a part operating through other channels. This
is a different question from the total effect of an adaptive policy. The
mediator's definition and observation time matter, and a missing engagement
record may itself be affected by contact intensity.

medflex supplies a natural-effect-model interface based on expanded data. Its
selected documentation explains that the original mediator is excluded from
the final natural-effect formula and that a full bootstrap refits working
models. Both details affect integration. Expanded records must retain their
link to the original account, and a bootstrap operating on those records
independently would not generally reproduce a bootstrap of the original
subjects and fitting procedure. Reusing a fixed first-stage fit would also
change the documented full-bootstrap operation. \cite{v12_sw_medflex}

The proposed adapter should therefore export an expansion identity, the
original sampling unit, the effect scale, the mediator definition and the
models refitted in each resample. Section~\ref{exec23:medflex-scope} traces
the expansion, full refitting and original-subject covariance, including the
cross-derivative contribution. That review does not make expanded rows new
independent information. The adapter should not market data expansion as a
solution to unmeasured mediator--outcome confounding. Sensitivity summaries
can accompany the decomposition, but they remain conditional analyses of
assumptions. A total-effect engine and a mediation engine can both be useful
without their outputs being interchangeable.

\section{What an R boundary must preserve}

An R process boundary is a plausible first integration route because it permits
evaluation of the existing estimator with its native dependencies. It is not
yet a selected or qualified deployment. The input should contain a stable
account identifier, the information and action clock, event/availability
times, the observation convention, the intervention specification and the
target's units. Serialization must distinguish absent values from absent
records and preserve factor levels, censoring codes and repeated-subject
structure.

The output needs more than an estimate. It must identify the exact target,
weights and truncations, nuisance-model choices, fitted sample, uncertainty
procedure and unavailable quantities. A bridge may fail before the estimator
runs, during fitting, or while translating a result. These failures require
different remedies. The first remedy is an environment or transport repair;
the second may require numerical diagnostics or revised assumptions; the
third requires correction of the adapter. None is a causal conclusion about
the underlying servicing policy.

Porting is an alternative only after the complete method boundary is clear.
Reimplementing a final regression while omitting its expansion, sequential
weighting, model-refitting bootstrap or stopping conditions would create a
different procedure. The advantage of a native port must be demonstrated on
the full workflow, including data transfer, model fitting, uncertainty and
failure handling. A faster final matrix multiplication is not enough.


